% !TeX root = ../main.tex
% =================================================================================
\chapter{Foundations}\label{chap:foundations}
% =================================================================================

This chapter introduces the concepts and technologies required for understanding this thesis.
\Cref{sec:foundations:first} introduces \ldots, and \cref{sec:foundations:second} covers \ldots

% Start every chapter with such an introductory paragraph that outlines its sections.
% A heading is never directly followed by another heading.

% ---------------------------------------------------------------------------------
\section{First Concept}\label{sec:foundations:first}
% ---------------------------------------------------------------------------------

Provide all (and exactly the) information that is necessary to understand the following chapters.
Assume as reader a fellow student of your degree program who did not study your specific subject.
Which concepts, terms, and definitions does this student have to know?
Which notations are standard in this area?
Only introduce what you actually need later, and cite the original sources.

For example, if your thesis builds on \glspl{llm}, briefly explain what they are and how they are used in \gls{ai}-assisted software development.
Acronyms are introduced on their first use (\gls{llm}) and abbreviated afterwards (\gls{llm}); the list of acronyms shows the page of their first use.

% ---------------------------------------------------------------------------------
\section{Second Concept}\label{sec:foundations:second}
% ---------------------------------------------------------------------------------

Introduce further foundations, e.g., technologies, frameworks, or existing tools that your work builds upon.
If your thesis extends an existing tool, describe it here in enough detail that the reader understands what you changed.
